% Introduction draft. Generated by make_paper_text.py -- do not edit numbers by hand.
\section{Introduction}
\label{sec:intro}

Retrieval-augmented generation (RAG) over document collections increasingly works on page images as well as on
extracted text. Visual retrievers such as ColPali~\cite{colpali} embed each page as a grid of patch vectors and
score queries by late interaction~\cite{colbert}, which handles tables, charts and layouts that optical character recognition (OCR) text loses.
The price is GPU compute and memory: every query is encoded and scored against every patch vector of every page, and
a vision-language generator then reads the page image. BM25~\cite{bm25} over OCR text answers a retrieval query on a
CPU in about a millisecond for our page pools of up to 1,286 pages, but needs OCR at indexing time
(0.26 GPU-s per page with docTR, more GPU time than ColPali indexing, which processes 16.2
pages/s). A system that
answers every question through the visual path, or through both, pays for the most expensive option on every query.

\subsection{Why routing in multimodal RAG is a supercomputing problem}
\label{sec:intro-hpc}

Three properties make serving multimodal RAG a high-performance computing problem rather than a modelling
detail; routing acts on the first and the third, index placement on the second.

\emph{Compute grows with the request rate.} Generation dominates the cost of every path. In a load test simulated
from our measured service times with dynamic batching, one GPU serving only the hybrid path cannot sustain
0.25\,queries/s at a p95 latency of 2\,s on 3 of 5 splits, and 100 queries/s needs 106 GPUs, 1.5$\times$
as many as content-aware routing.

\emph{Memory grows with the corpus.} ColPali stores 1,031 vectors per page; at 100k pages the index needs
24.6\,GiB in bf16, more than our 24\,GiB GPU, and a million pages (about 246\,GiB) exceed the
memory of most current accelerators. With exhaustive MaxSim, as measured here, such an index must be sharded, streamed from host memory or compressed;
hierarchical token pooling~\cite{tokenpooling}, evaluated here, shrinks it about fourfold with little loss of recall, and engines
such as PLAID~\cite{plaid} avoid exhaustive scoring (not evaluated here).
Measured on three older GPUs (Tesla P40) and two data-centre NVIDIA A10 GPUs, sharding keeps
exhaustive search at constant latency when the number of GPUs grows with the corpus (327\,ms for
50k pages on one Tesla P40, 328\,ms for 150k pages on three). Holding the index in host memory instead, two A10s searched
a million pages in batches of 64 at 0.44\,GPU-s per query, close to GPU-resident search
(94\% of its throughput where both were measured), which would need 16 A10s; at
6.4\,s per single query, this placement suits batch rather than interactive search. Routing
does not shrink the index, since our best router reads ColPali scores for every query; index placement is a separate
decision that the scaling study quantifies.

\emph{Latency targets constrain the design.} Interactive systems are judged on tail latency, so routing and
escalation should add as little serial work as possible. Our content-aware router adds 24.8\,ms of inference
before generation; we measure this overhead and charge it to every query.

Per-query routing is therefore a scheduling decision over GPU capacity: each query should get the cheapest context
that still answers it, and an operator should be able to trade quality for capacity at run time.

Cost-aware routing and cascades have been studied for large language models~\cite{frugalgpt,routellm} and for
adaptive retrieval~\cite{adaptiverag}; serving systems batch and schedule model inference under latency
targets~\cite{clipper,orca,vllm}, and efficient engines reduce late-interaction scoring cost~\cite{plaid}
(Section~\ref{sec:related}). We bring these threads together for multimodal document RAG, where the choice is
between context modalities that share one generator, and evaluate routing end to end on answer quality and
measured cost.
Routing, retrieval and generation are measured on a single GPU with Hugging Face (HF) Transformers generation, and capacity at
larger scale is simulated from measured service times and checked against live serving with vLLM on two A10 GPUs; multi-GPU sharding and host-memory placement of the index are
measured on two further nodes, with three older GPUs (Tesla P40) and with two NVIDIA A10 GPUs.

\subsection{Research questions}

\begin{description}
  \item[RQ1] What do the text, visual and hybrid paths cost, and how well do they answer, when measured end to end
  with real retrieval and the same generator?
  \item[RQ2] Can a router predict, before generation, which path answers a query well enough, and how close does
  it come to a per-query oracle?
  \item[RQ3] Can the cost--quality trade-off be controlled at run time, and can uncertain queries be escalated
  without adding serial latency?
  \item[RQ4] How does routing change the GPUs needed to meet a latency target, and how does the visual index scale
  with corpus size and the number of GPUs?
\end{description}

\subsection{Contributions}

\begin{itemize}
  \item An end-to-end evaluation that runs every path on every query of DocVQA~\cite{docvqa} and
  InfographicVQA~\cite{infovqa} with real BM25 and ColPali retrieval and a local generator, and measures answer
  quality as average normalised Levenshtein similarity (ANLS), GPU-seconds, energy and latency per query. All tables and figures are generated from cached outputs by
  released code.
  \item Outcome-based route labels, and the finding that keyword-derived labels agree with them no better than chance
  (Cohen's $\kappa = 0.01$).
  \item The Cost-Aware Visual Router (CAVR), a multi-task router that applies the standard cost-penalised decision rule
  $\arg\max_p \hat{q}_p - \lambda \bar{c}_p$ with $\lambda$ set at serving time, with its relation to
  budget-constrained routing made explicit; and CAVR-CE, a content-aware quality predictor built on a cross-encoder (CE) that uses 29\% less
  GPU compute than always-hybrid at a mean ANLS cost of 0.008, and 20\% less than a quality-blind mix
  of the visual and hybrid paths at equal ANLS; on a held-out confirmation set the saving over always-hybrid holds,
  but always-visual comes within 0.005 ANLS of the router at lower cost.
  \item An analysis of confidence-based escalation that separates pre-generation re-routing from serial and
  speculative cascades, with their latency and GPU cost computed from measured per-stage times.
  \item A serving study: capacity per GPU and GPUs needed at a p95 target under Poisson load (simulated from measured
  service times and checked against live vLLM serving on two A10 GPUs, where the policy comparison held but absolute
  capacity was up to 2.0$	imes$ lower than simulated), router throughput (measured), and visual-index latency up to 100k pages with GPU-resident,
  host-streamed and pooled placement (measured; larger indexes built from replicated page embeddings).
  \item Index placement beyond one GPU, measured: multi-GPU sharding on two nodes (up to three Tesla P40 and two NVIDIA A10 GPUs), with near-linear strong and weak scaling up to
  150k pages, and an index of a million pages held in host memory and streamed to two GPUs, with
  its latency and throughput compared with GPU-resident search.
\end{itemize}

We also report what does not work. With a shared generator the text path is not cheap; routing on the question text
alone is no better than always choosing the visual path; escalation on confidence signals reaches the same
operating point as lowering $\lambda$; and a per-query oracle reaches ANLS 0.630 at
0.486 GPU-s, far beyond any router we trained, because predicting per-query success remains hard.
A queue-driven controller for $\lambda$, evaluated in simulation under bursty load, gives only a small gain.
